\section{Peta Konsep Pemrograman Internet}

Mata kuliah Pemrograman Internet mencakup topik-topik utama yang saling terkait dan disusun mengikuti alur pembelajaran dari dasar ke lanjut \cite{mdn-learn}. Peta konsep ini membantu mahasiswa memvisualisasikan struktur materi secara keseluruhan dan memahami urutan belajar yang disarankan.

\begin{enumerate}
  \item \textbf{Bab II}: Landasan Teori --- Arsitektur web, HTTP, client-server, peran client-side dan server-side
  \item \textbf{Bab III}: HTML5 --- Struktur dokumen, elemen semantik, link, gambar, tabel
  \item \textbf{Bab IV}: HTML5 Form --- Form dan elemen input
  \item \textbf{Bab V}: CSS3 --- Selector, box model, warna, typography
  \item \textbf{Bab VI}: CSS3 Layout --- Flexbox, Grid, desain responsif
  \item \textbf{Bab VII}: JavaScript --- Sintaks, DOM, event handling
  \item \textbf{Bab VIII}: JavaScript --- Validasi form, integrasi dengan HTML/CSS
  \item \textbf{Bab IX}: Server-Side --- Pengenalan PHP, request-response, environment XAMPP
  \item \textbf{Bab X}: Basis Data --- Koneksi PHP-MySQL, operasi CRUD
  \item \textbf{Bab XI}: Session, Cookie, Autentikasi dasar
  \item \textbf{Bab XII}: Integrasi Aplikasi Web --- HTML/CSS/JS/PHP/MySQL
  \item \textbf{Bab XIII}: Keamanan Web --- Validasi input, prepared statement, sanitasi
  \item \textbf{Bab XIV}: Evaluasi dan Integrasi Kompetensi
\end{enumerate}

\textbf{Alur Pembelajaran:} Bab II--IV membangun fondasi arsitektur dan struktur konten; Bab V--VIII membangun antarmuka dan interaktivitas; Bab IX--XII membangun logika server, basis data, dan integrasi aplikasi; Bab XIII menambahkan lapisan keamanan; Bab XIV mengintegrasikan seluruh kompetensi melalui evaluasi komprehensif.

\begin{figure}[h]
\centering
\begin{tikzpicture}[node distance=1cm, font=\small]
  \node[client, draw] (n1) {Bab II--IV: Landasan + HTML};
  \node[client, draw, below=of n1] (n2) {Bab V--VIII: CSS + JavaScript};
  \node[server, draw, below=of n2] (n3) {Bab IX--XII: PHP + MySQL};
  \node[draw, rectangle, fill=orange!20, below=of n3] (n4) {Bab XIII: Keamanan};
  \node[draw, rectangle, fill=green!20, below=of n4] (n5) {Bab XIV: Evaluasi};
  \draw[arrow] (n1) -- (n2);
  \draw[arrow] (n2) -- (n3);
  \draw[arrow] (n3) -- (n4);
  \draw[arrow] (n4) -- (n5);
\end{tikzpicture}
\caption{Alur Pembelajaran Buku Ajar Pemrograman Internet}
\end{figure}
